\documentclass[11pt]{article}
\usepackage[a4paper,margin=1.15in]{geometry}
\usepackage{amsmath,amssymb,amsthm}
\usepackage[T1]{fontenc}
\usepackage{lmodern}
\usepackage[hidelinks]{hyperref}
\emergencystretch=3em

\newtheorem{theorem}{Theorem}
\newtheorem{proposition}{Proposition}
\newtheorem{lemma}{Lemma}
\newtheorem{definition}{Definition}
\newtheorem{remark}{Remark}
\newtheorem*{goal}{Statement}

% the rank-one meanings are the default; the rank-n section overrides and restores
\newcommand{\At}{A_t}
\newcommand{\Ft}{F_t}
\newcommand{\Nev}{\mathcal N^{\mathrm{ev}}}
\newcommand{\Uev}{U^{\mathrm{hyb,ev}}}
\newcommand{\HC}{\operatorname{HC}}
\newcommand{\Yc}{\check Y}
\newcommand{\Z}{\mathbb Z}
\newcommand{\ev}{\mathrm{ev}}
\newcommand{\qbinom}[2]{\genfrac{[}{]}{0pt}{}{#1}{#2}_{q}}
\newcommand{\Atev}{\mathbb A_t^{\mathrm{ev}}}
\newcommand{\Ftev}{\mathbb F_t^{\mathrm{ev}}}
\newcommand{\Nat}{\mathcal N_{\mathbb A_t}}
\newcommand{\Natev}{\mathcal N_{\mathbb A_t^{\mathrm{ev}}}}

\title{\textbf{The centre of the even hybrid family quantum $GL_n$}\\[2pt]
\large the statement for general $n$, and the case $n=1$ proved\thanks{The statement for
$n=1$ was fixed and signed before any proof was attempted, and the proof was checked against the
signed text; this document was assembled from those two records by the Almanac Editor of Project
Sandbox --- for the pilot edition, a software agent. The full record, including per-claim warrants
and machine verification, is the accompanying almanac edition.}}
\author{Tam\'as Hausel}
\date{}


\begin{document}

\maketitle



\begin{abstract}\noindent
A family quantum group carries, alongside the deformation parameter, independent toral parameters
$t_1,\dots,t_n$. Its \emph{even hybrid integral form} is the part of it that preserves every
integral Verma lattice; the question this paper is about is what the centre of that form looks like.
The answer proposed here is that the Harish--Chandra projection carries it isomorphically onto the
Weyl-invariant part of a \emph{Newton lattice} --- the functions on the multiplicative grid that take
integral values at every weight. Section 2 states this for $GL_n$. Sections 3 and 4 state and prove
it for $n=1$, where the root system is empty, the Weyl group is trivial, and the whole content is a
statement about integer-valued Laurent polynomials on a geometric grid: a bilateral form of the
quantum P\'olya theorem of Harman and Hopkins \cite{HarmanHopkins}. Section 5 is a short history.
\end{abstract}

\section{Introduction}

\noindent
Two things distinguish the family quantum group from the usual one. The toral parameters
$t_1,\dots,t_n$ are free, so the algebra is a family over a torus rather than a single object; and
one takes not the whole algebra but the \emph{hybrid} integral form --- the elements preserving
every integral Verma lattice --- and inside that, the torally even part. The centre of that form is
the object of study.

\medskip\noindent
The shape of the answer is the same in every rank, and it is worth saying before any notation. A
central element is determined by what it does to a highest-weight vector, weight by weight; that
assignment is a function on a grid of weights; and the condition of being \emph{integral} --- of
landing in $\At$ at every weight rather than merely in the fraction field --- cuts out a lattice of
such functions. \emph{The theorem is that the centre is exactly that lattice of integer-valued
functions, cut down by the Weyl symmetry.} Nothing is lost and nothing extra appears.

\medskip\noindent
The lattice in question has a long classical ancestry. That the integer-valued polynomials on
$\mathbb Z$ are free on the binomial coefficients is P\'olya's basis theorem
\cite[Prop.~1.1]{HarmanHopkins}; the standard reference for the subject is Cahen--Chabert
\cite{CahenChabert}. On a \emph{geometric} progression --- the grid that appears here --- the
corresponding basis is Gramain's \cite[Prop.~2.2]{Gramain}, for $q$ a natural number; the
$q$-analogue on the additive grid, with $q$ an indeterminate, is Harman--Hopkins
\cite{HarmanHopkins}. Section~5 places the present setting among them.

\medskip\noindent
Section~2 states the result for $GL_n$ in four clauses: a PBW basis, an intrinsic characterisation
of the integral form, the Harish--Chandra isomorphism onto the invariant Newton lattice, and a
coefficient-even refinement. \textbf{This general statement is not proved here.} It is the
programme's flagship, and it is stated so the rank-one case can be seen for what it is: not a
curiosity, but the first instance, with the root system empty.

\medskip\noindent
Sections~3 and~4 are the case $n=1$, stated and then proved in full. With $\Phi=\emptyset$ and
$W=\{1\}$ the algebra is a torus, the Harish--Chandra projection is the identity, and the three
clauses reduce to one substantial assertion: the even Newton lattice is spanned over $\At$ by the
shifted divided classes $\Yc^{u}\nu_r(\Yc)$.

\medskip\noindent
The rank-one case is small on purpose. It was carried from statement to proof to machine-checked
formalisation so that the cost and the failure modes of doing so could be measured; nine of the
twelve results in the accompanying edition carry a proof-kernel certificate. That apparatus is not
the subject of this paper and is not described here.

\section{The statement for $GL_n$}

\noindent
\emph{What follows is the programme's flagship statement. It is stated, not proved.} Its
individual ingredients are classical: the triangular decomposition and the freeness of the toral
part go back to Lusztig \cite[\S3.2]{Lusztig}; quantum Verma modules and their contravariant forms
are treated in De Concini--Procesi \cite[\S\S17.1--17.4]{DCP}; and over the fraction field the
Harish--Chandra map is an isomorphism onto the invariants by De Concini--Procesi
\cite[\S18.3]{DCP} and Jantzen \cite[Thm.~6.25 and \S6.26]{Jantzen} --- Jantzen's \S6.6 already
placing the image inside the \emph{even} toral invariants. Integral forms assembled, as here, from
a divided-power negative part and a toral lattice appear in Habiro--L\^e \cite[\S8G]{HabiroLe}.
What is new in the statement is the family parameter, the hybrid (Verma-preserving)
characterisation, and above all the assertion that over $\At$ the centre is the full
\emph{integer-valued} lattice --- not the monomial ring the generic statement would suggest.
\renewcommand{\At}{\mathbb A_t}\renewcommand{\Ft}{\mathbb F_t}

Fix $n\ge 1$ and put
\[
  q=v^{2},\qquad
  \At=\mathbb Z[v,v^{-1}][t_1^{\pm1},\ldots,t_n^{\pm1}],\qquad
  \Ft=\mathbb Q(v)[t_1^{\pm1},\ldots,t_n^{\pm1}],
\]
\[
  \Atev=\mathbb Z[q,q^{-1}][t_1^{\pm1},\ldots,t_n^{\pm1}],\qquad
  \Ftev=\mathbb Q(q)[t_1^{\pm1},\ldots,t_n^{\pm1}],
\]
with $v$ an indeterminate and $t_1,\ldots,t_n$ algebraically independent.
Let $\varepsilon_1,\ldots,\varepsilon_n$ be the standard basis of
$\mathbb Z^{n}$, equipped with $(\varepsilon_a,\varepsilon_b)=\delta_{ab}$;
put $\alpha_i=\varepsilon_i-\varepsilon_{i+1}$ ($1\le i<n$) and
$\beta_{ab}=\varepsilon_a-\varepsilon_b$ ($1\le a<b\le n$), and let
$\Phi^{+}=\{\beta_{ab}:1\le a<b\le n\}\subset\mathbb Z^{n}$ be the set of
positive roots, listed from left to right in the order
\[
  (1,2),\,(1,3),\ldots,(1,n),\,(2,3),\ldots,(n-1,n).
\]
Arrays are written $\mathbf a=(a_{ab})_{a<b}\in\mathbb Z_{\ge0}^{\Phi^{+}}$,
with $\mathbf 0$ the zero array.  For $Z$ the centre of an algebra we write
$Z(-)$.  Quantum integers and factorials are
\[
  [r]_v=\frac{v^{r}-v^{-r}}{v-v^{-1}},\qquad
  [r]_v!=\prod_{j=1}^{r}[j]_v,\qquad
  [0]_v!=1 .
\]

\begin{definition}[Reduced double of family quantum $GL_n$]\label{def:utilde}
$\widetilde U_{\Ft}$ is the $\Ft$-algebra with invertible pairwise commuting
generators $K_a^{+},K_a^{-}$ ($1\le a\le n$) and generators $e_i,F_i$
($1\le i<n$), subject to
\[
  K_a^{+}K_a^{-}=t_a,\qquad
  K_a^{\pm}\,e_i\,(K_a^{\pm})^{-1}=v^{\pm(\varepsilon_a,\alpha_i)}e_i,\qquad
  K_a^{\pm}\,F_i\,(K_a^{\pm})^{-1}=v^{\mp(\varepsilon_a,\alpha_i)}F_i,
\]
\[
  e_iF_j-F_je_i=\delta_{ij}\bigl(K_i^{+}K_{i+1}^{-}-K_i^{-}K_{i+1}^{+}\bigr),
\]
\[
  e_ie_j=e_je_i,\qquad F_iF_j=F_jF_i\qquad(|i-j|>1),
\]
\[
  e_i^{2}e_j-(v+v^{-1})e_ie_je_i+e_je_i^{2}=0,\qquad
  F_i^{2}F_j-(v+v^{-1})F_iF_jF_i+F_jF_i^{2}=0\qquad(|i-j|=1).
\]
Define interval root vectors by $e_{a,a+1}=e_a$, $F_{a,a+1}=F_a$ and, for
$b>a+1$,
\[
  e_{ab}=\frac{e_{a,b-1}\,e_{b-1,b}-v^{-1}e_{b-1,b}\,e_{a,b-1}}{v-v^{-1}},
  \qquad
  F_{ab}=F_{b-1,b}\,F_{a,b-1}-v\,F_{a,b-1}\,F_{b-1,b}.
\]
For $\mathbf a\in\mathbb Z_{\ge0}^{\Phi^{+}}$ the divided and ordinary PBW
words are
\[
  F^{(\mathbf a)}=\prod_{a<b}^{\longleftarrow}\frac{F_{ab}^{\,a_{ab}}}{[a_{ab}]_v!},
  \qquad
  e^{\mathbf a}=\prod_{a<b}^{\longrightarrow}e_{ab}^{\,a_{ab}},
\]
where the arrows read the displayed root list backwards and forwards.  Put
$\nu(\mathbf a)=\sum_{a<b}a_{ab}\,\beta_{ab}\in\mathbb Z^{n}$, and for
$\mu=\sum_a\mu_a\varepsilon_a\in\mathbb Z^{n}$ put
$K^{+}_{\mu}=\prod_a (K_a^{+})^{\mu_a}$.
\end{definition}

\begin{definition}[Torally even part]\label{def:even}
$\widetilde U_{\Ft}$ carries the $\mathbb Z^{n}/2\mathbb Z^{n}$-grading
determined on generators by
\[
  \deg K_a^{+}=\varepsilon_a,\qquad
  \deg K_a^{-}=-\varepsilon_a,\qquad
  \deg K^{+}_{\mu}=\mu,\qquad
  \deg e_i=0,\qquad
  \deg F_i=-\alpha_i,\qquad
  \deg t_a=0;
\]
every defining relation is homogeneous modulo $2\mathbb Z^{n}$ (in the mixed
relation both toral terms have degree $-\alpha_i\equiv\alpha_i$, the degree
of its left side), so the grading is well defined, and the degree-zero part
is a subalgebra.  The interval vectors are homogeneous with
$\deg e_{ab}=0$ and $\deg F_{ab}=-\beta_{ab}$, so
$\deg F^{(\mathbf a)}=-\nu(\mathbf a)$.  $U_{\Ft}$ is the degree-zero part
of $\widetilde U_{\Ft}$.
\end{definition}

\begin{definition}[Newton lattice; the hybrid form]\label{def:newton}
Put
\[
  Y_a=q^{\,n-a}\,t_a^{-1}\,(K_a^{+})^{2}\in\widetilde U_{\Ft},
  \qquad\text{equivalently}\qquad
  (K_a^{+})^{2}=q^{\,a-n}\,t_a\,Y_a .
\]
Granting the basis in clause (i) below, the $Y_a$ are algebraically
independent invertible elements of $U_{\Ft}$, so
$\Ft[Y_1^{\pm1},\ldots,Y_n^{\pm1}]$ is a Laurent polynomial ring; for
$\lambda\in\mathbb Z^{n}$ let
\[
  \ev_\lambda:\Ft[Y_1^{\pm1},\ldots,Y_n^{\pm1}]\longrightarrow\Ft,
  \qquad
  \ev_\lambda(Y_a)=q^{\lambda_a+n-a}\,t_a^{-1},
\]
be the $\Ft$-algebra map.  The \emph{Newton lattice} is
\[
  \Nat=\bigl\{\,f\in\Ft[Y_1^{\pm1},\ldots,Y_n^{\pm1}]:
  \ev_\lambda(f)\in\At\ \text{for every }\lambda\in\mathbb Z^{n}\,\bigr\},
\]
and the \emph{even hybrid form} is the $\At$-submodule
\[
  U_{\At}\;=\;\sum_{\mathbf a,\mathbf b\in\mathbb Z_{\ge0}^{\Phi^{+}}}
  F^{(\mathbf a)}\,K^{+}_{\nu(\mathbf a)}\,\Nat\,e^{\mathbf b}
  \;\subset\;U_{\Ft},
\]
each summand torally even since
$\deg\bigl(F^{(\mathbf a)}K^{+}_{\nu(\mathbf a)}\bigr)=0$; only finite sums
occur.  The symmetric group $S_n$ fixes $\Ft$ and permutes the variables
$Y_1,\ldots,Y_n$; the invariant Newton lattice is
$(\Nat)^{S_n}=\{f\in\Nat: wf=f\ \text{for every }w\in S_n\}$.
\end{definition}

\begin{definition}[Integral Verma lattices; Harish--Chandra
projection]\label{def:verma-hc}
Let $\widetilde U_{\Ft}^{\ge0}\subset\widetilde U_{\Ft}$ be the subalgebra
generated by all $K_a^{\pm}$ and all $e_i$.  For $\lambda\in\mathbb Z^{n}$
let $\chi_\lambda:\widetilde U_{\Ft}^{\ge0}\to\Ft$ be the $\Ft$-algebra map
with
\[
  \chi_\lambda(K_a^{+})=v^{\lambda_a},\qquad
  \chi_\lambda(K_a^{-})=t_a\,v^{-\lambda_a},\qquad
  \chi_\lambda(e_i)=0,
\]
writing $\Ft^{\chi_\lambda}$ for $\Ft$ viewed as a
$\widetilde U_{\Ft}^{\ge0}$-module through $\chi_\lambda$, let
$M_{\Ft}(\lambda)
=\widetilde U_{\Ft}\otimes_{\smash{\widetilde U_{\Ft}^{\ge0}}}\Ft^{\chi_\lambda}$
be the induced module, with highest vector $m_\lambda=1\otimes1$; its
integral lattice is
\[
  M_{\At}(\lambda)\;=\;\sum_{\mathbf a\in\mathbb Z_{\ge0}^{\Phi^{+}}}
  \At\,F^{(\mathbf a)}\,K^{+}_{\nu(\mathbf a)}\,m_\lambda .
\]
Granting the basis in clause (i) below, every element of $U_{\Ft}$ has a
unique PBW coefficient on the row $\mathbf a=\mathbf b=\mathbf 0$, and that
coefficient lies in $\Ft[Y_1^{\pm1},\ldots,Y_n^{\pm1}]$; the
\emph{Harish--Chandra projection} is the $\Ft$-linear map
\[
  \HC:U_{\Ft}\longrightarrow\Ft[Y_1^{\pm1},\ldots,Y_n^{\pm1}]
\]
retaining that coefficient.
\end{definition}

\begin{goal}\begin{enumerate}
  \item[(i)] The PBW sums are direct: the monomials
  $F^{(\mathbf a)}(K_1^{+})^{\eta_1}\cdots(K_n^{+})^{\eta_n}e^{\mathbf b}$,
  for $\mathbf a,\mathbf b\in\mathbb Z_{\ge0}^{\Phi^{+}}$ and
  $\eta\in\mathbb Z^{n}$, form an $\Ft$-basis of $\widetilde U_{\Ft}$; the
  degree-zero part decomposes as
  \[
    U_{\Ft}\;=\;\bigoplus_{\mathbf a,\mathbf b\in\mathbb Z_{\ge0}^{\Phi^{+}}}
    F^{(\mathbf a)}\,K^{+}_{\nu(\mathbf a)}\,
    \Ft[Y_1^{\pm1},\ldots,Y_n^{\pm1}]\,e^{\mathbf b};
  \]
  the sum defining $U_{\At}$ is direct, $U_{\At}$ is an $\At$-subalgebra of
  $U_{\Ft}$ with $\Ft\otimes_{\At}U_{\At}\simeq U_{\Ft}$, and each sum
  defining $M_{\At}(\lambda)$ is direct.
  \item[(ii)] The even hybrid form is the common preserver of the integral
  Verma lattices:
  \[
    U_{\At}\;=\;\bigl\{\,x\in U_{\Ft}:
    x\,M_{\At}(\lambda)\subseteq M_{\At}(\lambda)\ \text{for every }
    \lambda\in\mathbb Z^{n}\,\bigr\}.
  \]
  \item[(iii)] The intrinsic centre satisfies
  $Z(U_{\At})=Z(\widetilde U_{\Ft})\cap U_{\At}$, and Harish--Chandra
  projection restricts to an isomorphism of $\At$-algebras
  \[
    \HC:\;Z\bigl(U_{\At}\bigr)\;\xrightarrow{\ \sim\ }\;
    \bigl(\Nat\bigr)^{S_n}.
  \]
  \item[(iv)] With
  \[
    \Natev=\bigl\{\,f\in\Ftev[Y_1^{\pm1},\ldots,Y_n^{\pm1}]:
    \ev_\lambda(f)\in\Atev\ \text{for every }\lambda\in\mathbb Z^{n}\,\bigr\}
  \]
  and $\mathfrak Z_{\Atev}=\{C\in Z(U_{\At}):\HC(C)\in(\Natev)^{S_n}\}$ the
  coefficient-even central form, $\HC$ restricts to an isomorphism of
  $\Atev$-algebras
  $\mathfrak Z_{\Atev}\xrightarrow{\ \sim\ }(\Natev)^{S_n}$, and
  multiplication induces
  $\At\otimes_{\Atev}\mathfrak Z_{\Atev}\xrightarrow{\ \sim\ }Z(U_{\At})$.
\end{enumerate}
\end{goal}

\renewcommand{\At}{A_t}\renewcommand{\Ft}{F_t}

\section{The case $n=1$}

\noindent
Here $\Phi=\emptyset$, $W=\{1\}$, and the cocharacter lattice is $\Lambda=\mathbb Z$. The four
clauses above collapse: there are no root vectors, so the PBW basis is the toral monomials alone;
the Harish--Chandra projection is the identity; and the Weyl invariance is vacuous. What remains is
clause~(iii), and it is not vacuous at all.

Let $G=GL_1$: the root system is $\Phi=\emptyset$, the Weyl group is
$W=\{1\}$, the cocharacter lattice is $\Lambda=\mathbb Z$, and
$X_\Lambda=\mathbb Z$. Let
\[
  \At=\mathbb Z[v^{\pm1},t^{\pm1}],\qquad
  \Ft=\mathbb Q(v)[t^{\pm1}],\qquad
  q=v^{2},
\]
with $v,t$ independent indeterminates.

\begin{definition}[Hybrid family quantum $GL_1$]\label{def:U}
$U$ is the $\Ft$-algebra with invertible pairwise commuting generators
$K_{n,+},K_{n,-}$ ($n\in\mathbb Z$) subject to
\[
  K_{0,\pm}=1,\qquad
  K_{m+n,\pm}=K_{m,\pm}K_{n,\pm},\qquad
  K_{n,+}K_{n,-}=t^{\,n}.
\]
There are no generators $e_i,F_i$, since $\Phi=\emptyset$. Writing
$T^{\pm}:=K_{1,\pm}$, the monomials $(T^{+})^{w}$ ($w\in\mathbb Z$) form a
$\Ft$-basis of $U$.
\end{definition}

\begin{definition}[Even Newton lattice]\label{def:nev}
A \emph{toral} element is an element of $\Ft[(T^{+})^{\pm1}]$; it is \emph{even}
when all its toral exponents are even --- equivalently, it is a Laurent
polynomial in $\Yc:=(T^{+})^{2}$. For $\chi\in X_\Lambda$ let
$\mathrm{ev}_\chi$ be the $\At$-linear map on even toral elements determined by
$\mathrm{ev}_\chi(\Yc)=q^{\,\chi}$. The \emph{even Newton lattice} is the
$\At$-submodule
\[
  \Nev=\bigl\{\,f\in\Ft[\Yc^{\pm1}]\;:\;
  \mathrm{ev}_\chi(f)\in\At\ \text{ for every }\chi\in\mathbb Z\,\bigr\}.
\]
\end{definition}

\begin{definition}[Even hybrid family integral form]\label{def:uev}
Since $\Phi^{+}=\emptyset$, the \emph{even hybrid family integral form} is the
$\At$-submodule
\[
  \Uev\;=\;\Nev\;\subset\;U,
\]
and its \emph{central part} is
$Z(\Uev)=\Uev\cap Z\bigl(U\otimes_{\Ft}\operatorname{Frac}(\At)\bigr)$.
\end{definition}

\begin{definition}[Harish--Chandra projection; shifted Weyl action]\label{def:hc}
Every element of $U$ is toral, so the \emph{Harish--Chandra projection} $\HC$ is
the identity map on $U$. With
\[
  Y_n=t^{-n}K_{n,+}^{2}\qquad(n\in\mathbb Z),
\]
every even toral element is a Laurent polynomial in the $Y_n$, and the
\emph{shifted Weyl action} of $W$ is the $\Ft$-algebra action
$w(Y_n)=Y_{wn}$; set $(\Nev)^{W}=\{f\in\Nev: wf=f\ \text{for all}\ w\in W\}$.
\end{definition}

\begin{definition}[Newton divided classes]\label{def:nu}
For $r\ge0$ set
\[
  \nu_{r}(X)=\prod_{s=0}^{r-1}\frac{X-q^{s}}{q^{r}-q^{s}}\;\in\;\mathbb Q(v)[X],
  \qquad \nu_0(X)=1 .
\]
\end{definition}

\begin{goal}\label{goal:main}
For $G=GL_1$, with $\Lambda=\mathbb Z$:
\begin{enumerate}
  \item[(i)] $\Uev$ is an $\At$-subalgebra of $U$;
  \item[(ii)] the Harish--Chandra projection restricts to an isomorphism of
  $\At$-algebras
  \[
    \HC:\;Z\bigl(\Uev\bigr)\;\xrightarrow{\ \sim\ }\;\bigl(\Nev\bigr)^{W};
  \]
  \item[(iii)] $\displaystyle
    \Nev=\sum_{u\in\mathbb Z}\ \sum_{r\ge0}\ \At\;\Yc^{\,u}\,\nu_{r}(\Yc)$.
\end{enumerate}
\end{goal}

\section{Proof of the case $n=1$}

Throughout, $v$ and $t$ are independent indeterminates and
\[
  \At=\Z[v^{\pm1},t^{\pm1}],\qquad
  \Ft=\mathbb Q(v)[t^{\pm1}],\qquad
  q=v^{2}.
\]
Thus $\At\subseteq\Ft$ is a subring, and $\Z[q^{\pm1}]=\Z[v^{\pm2}]\subseteq
\Z[v^{\pm1}]\subseteq\At$.

\emph{The algebra $U$.}\;
$U$ is the $\Ft$-algebra with invertible pairwise-commuting generators
$K_{n,+},K_{n,-}$ ($n\in\Z$) subject to
$K_{0,\pm}=1$, $K_{m+n,\pm}=K_{m,\pm}K_{n,\pm}$ and $K_{n,+}K_{n,-}=t^{\,n}$.
Writing $T^{\pm}=K_{1,\pm}$, the monomials $(T^{+})^{w}$ ($w\in\Z$) form a
$\Ft$-basis; $U$ is commutative. A \emph{toral} element is an element of
$\Ft[(T^{+})^{\pm1}]$; it is \emph{even} when it is a Laurent polynomial in
\[
  \Yc:=(T^{+})^{2}.
\]
Even toral elements thus form the Laurent-polynomial algebra
$\Ft[\Yc^{\pm1}]$.

\emph{Evaluations.}\;
For $\chi\in\Z$, let $\ev_\chi\colon\Ft[\Yc^{\pm1}]\to\Ft$ be the
$\Ft$-algebra homomorphism substituting the unit $q^{\chi}\in\At$ for $\Yc$,
i.e.
\[
  \ev_\chi\Bigl(\textstyle\sum_{u}c_u\Yc^{u}\Bigr)=\sum_u c_u\,q^{u\chi}
  \qquad(c_u\in\Ft).
\]
It is the unique $\At$-linear ring homomorphism with $\ev_\chi(\Yc)=q^{\chi}$,
and it fixes $\Ft$ (hence $\At$) pointwise. The \emph{even Newton lattice} is
\[
  \Nev=\bigl\{\,f\in\Ft[\Yc^{\pm1}]:\ \ev_\chi(f)\in\At\ \text{for every }
  \chi\in\Z\,\bigr\}.
\]
Because $\Phi^{+}=\emptyset$ at $GL_1$, the \emph{even hybrid family integral
form} is $\Uev=\Nev\subseteq U$, with central part
$Z(\Uev)=\Uev\cap Z\bigl(U\otimes_{\Ft}\operatorname{Frac}(\At)\bigr)$. Every
element of $U$ is toral, so the Harish--Chandra projection $\HC$ is the identity
map on $U$; the Weyl group is $W=\{1\}$, acting trivially, so
$(\Nev)^{W}=\Nev$.

\emph{Newton divided classes.}\;
For $r\ge0$,
\[
  \nu_{r}(X)=\prod_{s=0}^{r-1}\frac{X-q^{s}}{q^{r}-q^{s}}\ \in\ \mathbb Q(v)[X],
  \qquad \nu_0(X)=1 .
\]
Each $\nu_r$ has degree exactly $r$ (its leading coefficient
$\prod_{s=0}^{r-1}(q^{r}-q^{s})^{-1}$ is a nonzero element of $\mathbb Q(v)$),
so $\{\nu_0,\dots,\nu_d\}$ is a $\Ft$-basis of the module
$\Ft[\Yc]_{\le d}$ of even toral polynomials of degree $\le d$.

\medskip
We write $\Z[q^{\pm1}]$-\emph{integral} for ``lies in $\Z[q^{\pm1}]$''.
For $m\in\Z$ and $r\ge0$ put
\[
  B_r(m):=\nu_r(q^{m})=\prod_{s=0}^{r-1}\frac{q^{m}-q^{s}}{q^{r}-q^{s}},
  \qquad B_0(m)=1 .
\]

\subsection{Clause (i): $\Uev$ is an $\At$-subalgebra}

\begin{proposition}
$\Uev=\Nev$ is an $\At$-subalgebra of $U$. (This is the bilateral,
$t$-spectator analogue of the fact that quantum integer-valued polynomials
form a ring \cite[\S1]{HarmanHopkins}.)
\end{proposition}

\begin{proof}
For each $\chi\in\Z$ the map $\ev_\chi\colon\Ft[\Yc^{\pm1}]\to\Ft$ is a ring
homomorphism, and $\At\subseteq\Ft$ is a subring; hence the preimage
$\ev_\chi^{-1}(\At)$ is a subring of $\Ft[\Yc^{\pm1}]$. It contains $\At$,
because $\ev_\chi$ fixes $\At$ pointwise, so $\ev_\chi^{-1}(\At)$ is an
$\At$-subalgebra. By definition
\[
  \Nev=\bigcap_{\chi\in\Z}\ev_\chi^{-1}(\At)
\]
is an intersection of $\At$-subalgebras of $\Ft[\Yc^{\pm1}]$, hence itself an
$\At$-subalgebra (it contains $1$ since $\ev_\chi(1)=1\in\At$, and is closed
under sums, products and $\At$-scalars). As $\Ft[\Yc^{\pm1}]\subseteq U$, this
is an $\At$-subalgebra of $U$.
\end{proof}

\subsection{Clause (ii): the Harish--Chandra isomorphism is the identity}

\begin{proposition}
The Harish--Chandra projection restricts to an isomorphism of $\At$-algebras
$\HC\colon Z(\Uev)\xrightarrow{\ \sim\ }(\Nev)^{W}$, and this isomorphism is the
identity map of $\Nev$. (The rank-one, $\Phi^{+}=\emptyset$ degeneration of the
even hybrid Harish--Chandra isomorphism established at $GL_2$ in
\cite{SandboxGL2}; the classical ancestor is the Harish--Chandra isomorphism
itself, whose quantum-group form the $GL_2$ record cites.)
\end{proposition}

\begin{proof}
The algebra $U$ is commutative, being spanned by monomials in the pairwise
commuting invertible generators $K_{n,\pm}$. Base change is compatible with
commutativity, so $U\otimes_{\Ft}\operatorname{Frac}(\At)$ is a commutative ring
and equals its own centre. Therefore
\[
  Z(\Uev)=\Uev\cap Z\bigl(U\otimes_{\Ft}\operatorname{Frac}(\At)\bigr)
        =\Uev\cap\bigl(U\otimes_{\Ft}\operatorname{Frac}(\At)\bigr)
        =\Uev=\Nev,
\]
the middle equality because $U\hookrightarrow U\otimes_{\Ft}\operatorname{Frac}(\At)$
is injective --- $\Ft\hookrightarrow\operatorname{Frac}(\At)=\mathbb Q(v,t)$ and $U$ is
free over $\Ft$ --- so $\Uev=\Nev\subseteq U$ is literally a subset of the base-changed
algebra.
Every element of $U$ is toral, so $\HC=\mathrm{id}_U$, and its restriction to
$Z(\Uev)=\Nev$ is $\mathrm{id}_{\Nev}$. Finally $W=\{1\}$ acts trivially, so
$(\Nev)^{W}=\Nev$. Hence
\[
  \HC\colon Z(\Uev)=\Nev\ \xrightarrow{\ \mathrm{id}\ }\ \Nev=(\Nev)^{W}
\]
is the identity map of the $\At$-algebra $\Nev$ (an $\At$-algebra by clause~(i)),
which is manifestly an isomorphism of $\At$-algebras.
\end{proof}

\noindent
There is no content here beyond commutativity: at $GL_1$ the rational centre is
everything, $\HC$ is the identity, and Weyl invariance is vacuous.

\subsection{Clause (iii): the Newton presentation}

Write
\[
  S\ :=\ \sum_{u\in\Z}\ \sum_{r\ge0}\ \At\,\Yc^{\,u}\,\nu_{r}(\Yc)\ \subseteq\
  \Ft[\Yc^{\pm1}]
\]
for the $\At$-submodule generated by the shifted divided classes
$\Yc^{u}\nu_r(\Yc)$.

\begin{theorem}[Even Newton presentation for $GL_1$]\label{thm:iii}
$\displaystyle \Nev=S=\sum_{u\in\Z}\sum_{r\ge0}\At\,\Yc^{\,u}\,\nu_r(\Yc).$
(The bilateral quantum P\'olya theorem of Harman--Hopkins
\cite[\S\S1,4]{HarmanHopkins}, transported to the multiplicative grid; the
precise dictionary is the Attribution remark below.)
\end{theorem}

The proof occupies the rest of this section: the grid lemma, the forward
inclusion $S\subseteq\Nev$, the windowed reverse inclusion for polynomials, and
the Laurent-clearing step.

\subsubsection{The grid values are $\Z[q^{\pm1}]$-integral}

\begin{lemma}\label{lem:grid}
For all $m\in\Z$ and $r\ge0$ one has $B_r(m)=\nu_r(q^{m})\in\Z[q^{\pm1}]$
(the Gaussian-binomial value pattern of \cite[\S1]{HarmanHopkins}, extended
to negative nodes as their \S4 bilateral localization requires).
Explicitly:
\begin{enumerate}
\item[(a)] if $0\le m<r$ then $B_r(m)=0$;
\item[(b)] if $m\ge r\ge0$ then $B_r(m)=\qbinom{m}{r}\in\Z[q]$;
\item[(c)] if $m=-n$ with $n\ge1$ then
  \[
    B_r(-n)=(-1)^{r}\,q^{-nr-\binom{r}{2}}\,\qbinom{n+r-1}{r}\ \in\ \Z[q^{\pm1}].
  \]
\end{enumerate}
\end{lemma}

\begin{proof}
\emph{(a)}\; If $0\le m<r$, the factor of the product at $s=m$ has numerator
$q^{m}-q^{m}=0$, so $B_r(m)=0$.

\emph{(b)}\; Assume $m\ge r\ge0$. Pulling $q^{s}$ out of each numerator and
denominator factor,
\[
  B_r(m)=\prod_{s=0}^{r-1}\frac{q^{s}(q^{m-s}-1)}{q^{s}(q^{r-s}-1)}
        =\prod_{s=0}^{r-1}\frac{q^{m-s}-1}{q^{r-s}-1}.
\]
Re-indexing by $i=r-s$ (so $i$ runs from $1$ to $r$),
\[
  B_r(m)=\prod_{i=1}^{r}\frac{q^{(m-r)+i}-1}{q^{i}-1}=\qbinom{m}{r},
\]
the Gaussian binomial coefficient. It lies in $\Z[q]$: this is standard from the
$q$-Pascal recurrence $\qbinom{m}{r}=\qbinom{m-1}{r-1}+q^{r}\qbinom{m-1}{r}$ with
$\qbinom{m}{0}=1$, by induction on $m$.

\emph{(c)}\; Let $m=-n$, $n\ge1$. In each factor write the numerator as
$q^{-n}-q^{s}=-q^{-n}(q^{n+s}-1)$ and the denominator as
$q^{r}-q^{s}=q^{s}(q^{r-s}-1)$. Then
\[
  B_r(-n)=\prod_{s=0}^{r-1}\frac{-q^{-n}(q^{n+s}-1)}{q^{s}(q^{r-s}-1)}
        =(-1)^{r}\,q^{-nr}\,q^{-\sum_{s=0}^{r-1}s}
          \prod_{s=0}^{r-1}\frac{q^{n+s}-1}{q^{r-s}-1}.
\]
Here $\sum_{s=0}^{r-1}s=\binom{r}{2}$. For the remaining product, the
denominators give $\prod_{s=0}^{r-1}(q^{r-s}-1)=\prod_{i=1}^{r}(q^{i}-1)$ and the
numerators give $\prod_{s=0}^{r-1}(q^{n+s}-1)=\prod_{j=n}^{n+r-1}(q^{j}-1)$, so
\[
  \prod_{s=0}^{r-1}\frac{q^{n+s}-1}{q^{r-s}-1}
  =\frac{\prod_{j=n}^{n+r-1}(q^{j}-1)}{\prod_{i=1}^{r}(q^{i}-1)}
  =\frac{[n+r-1]_q!\,/\,[n-1]_q!}{[r]_q!}=\qbinom{n+r-1}{r},
\]
where $[k]_q=(q^{k}-1)/(q-1)$ and the common powers of $(q-1)$ cancel. The
Gaussian binomial $\qbinom{n+r-1}{r}\in\Z[q]$ by part (b). Collecting the
prefactor gives
\[
  B_r(-n)=(-1)^{r}\,q^{-nr-\binom{r}{2}}\,\qbinom{n+r-1}{r}\in\Z[q^{\pm1}].
  \qedhere
\]
\end{proof}

\subsubsection{Forward inclusion $S\subseteq\Nev$}

\begin{lemma}\label{lem:fwd}
For every $u\in\Z$ and $r\ge0$, $\Yc^{u}\nu_r(\Yc)\in\Nev$. Hence $S\subseteq\Nev$.
(The transported forward half of \cite[\S1]{HarmanHopkins}: $q$-binomials are
quantum integer-valued.)
\end{lemma}

\begin{proof}
Fix $u,r$ and $\chi\in\Z$. Since $\ev_\chi$ is a ring homomorphism with
$\ev_\chi(\Yc)=q^{\chi}$,
\[
  \ev_\chi\bigl(\Yc^{u}\nu_r(\Yc)\bigr)=q^{u\chi}\,\nu_r(q^{\chi})
  =q^{u\chi}\,B_r(\chi).
\]
By Lemma~\ref{lem:grid}, $B_r(\chi)\in\Z[q^{\pm1}]\subseteq\At$, and $q^{u\chi}$
is a unit of $\At$; so the product lies in $\At$. As $\chi$ was arbitrary,
$\Yc^{u}\nu_r(\Yc)\in\Nev$. The lattice $\Nev$ is an $\At$-module (indeed an
$\At$-algebra, clause~(i)), so it contains every $\At$-combination of these
generators, i.e.\ $S\subseteq\Nev$.
\end{proof}

\subsubsection{Reverse inclusion for polynomials: windowed Newton interpolation}

\begin{lemma}\label{lem:rev}
Let $g\in\Nev$ be an even toral \emph{polynomial}, $g\in\Ft[\Yc]$, of degree
$\le d$. Then $g=\sum_{r=0}^{d}c_r\,\nu_r(\Yc)$ with all $c_r\in\At$. In
particular $g\in\sum_{r\ge0}\At\,\nu_r(\Yc)\subseteq S$.
(P\'olya's windowed-interpolation argument as $q$-deformed in the proof of
\cite[\S1]{HarmanHopkins}, run on the multiplicative grid with $t$ a
spectator.)
\end{lemma}

\begin{proof}
Because $\{\nu_0,\dots,\nu_d\}$ is a $\Ft$-basis of $\Ft[\Yc]_{\le d}$, there are
unique $c_0,\dots,c_d\in\Ft$ with
\[
  g=\sum_{r=0}^{d}c_r\,\nu_r(\Yc).
\]
Evaluate on the window $\chi=m$ for $m=0,1,\dots,d$. As $\ev_m$ is
$\Ft$-linear and $\ev_m(\nu_r(\Yc))=B_r(m)$,
\[
  \ev_m(g)=\sum_{r=0}^{d}B_r(m)\,c_r
  =\sum_{r=0}^{d}V[m,r]\,c_r,
  \qquad V[m,r]:=B_r(m)\ \ (0\le m,r\le d).
\]
By Lemma~\ref{lem:grid}, every entry $V[m,r]\in\Z[q^{\pm1}]$; moreover $V$ is
\emph{lower unitriangular}: $V[m,r]=B_r(m)=0$ for $m<r$ (part (a)), and
$V[m,m]=B_m(m)=\prod_{s=0}^{m-1}\frac{q^{m}-q^{s}}{q^{m}-q^{s}}=1$. Hence
$\det V=1$, and the inverse $V^{-1}$ is again lower unitriangular with entries in
$\Z[q^{\pm1}]$ (back-substitution never leaves $\Z[q^{\pm1}]$, and by uniqueness
this inverse coincides with the inverse taken over the larger ring $\Ft$).

Writing $\mathbf c=(c_0,\dots,c_d)^{\mathsf T}$ and
$\mathbf v=(\ev_0(g),\dots,\ev_d(g))^{\mathsf T}$, we have $\mathbf v=V\mathbf c$
in $\Ft^{d+1}$, hence
\[
  \mathbf c=V^{-1}\mathbf v .
\]
Since $g\in\Nev$, each $\ev_m(g)\in\At$, so $\mathbf v\in\At^{d+1}$; and
$V^{-1}$ has entries in $\Z[q^{\pm1}]\subseteq\At$. Therefore
$\mathbf c=V^{-1}\mathbf v\in\At^{d+1}$, i.e.\ every $c_r\in\At$.

\smallskip
\emph{The $t$-spectator argument.}\;
The map $\ev_m$ substitutes only $\Yc\mapsto q^{m}=v^{2m}$; it does not touch $v$
or $t$. Thus $t$ never enters the matrix $V$, whose entries lie in
$\Z[q^{\pm1}]$. Concretely, expand each coordinate along the free
$\Z[v^{\pm1}]$-basis of $t$-powers: $c_r=\sum_{k}c_{r,k}\,t^{k}$ with
$c_{r,k}\in\mathbb Q(v)$, so that
$\ev_m(g)=\sum_{k}\bigl(\sum_{r}B_r(m)\,c_{r,k}\bigr)t^{k}$. Membership
$\ev_m(g)\in\At=\Z[v^{\pm1},t^{\pm1}]$ means that for each fixed $k$ the
$t^{k}$-coefficient $\sum_r B_r(m)c_{r,k}$ lies in $\Z[v^{\pm1}]$ for all
$m\in\{0,\dots,d\}$. For each fixed $k$ this is precisely the unitriangular
system above, over $\Z[q^{\pm1}]$ with right-hand sides in $\Z[v^{\pm1}]$; it
gives $c_{r,k}\in\Z[v^{\pm1}]$. Hence
$c_r=\sum_k c_{r,k}t^{k}\in\Z[v^{\pm1},t^{\pm1}]=\At$, as claimed. This is the
same conclusion reached above; we spell out the coefficient-wise reading to make
the spectator role of $t$ explicit.

\smallskip
\emph{Base change $\Z[q^{\pm1}]\to\Z[v^{\pm1}]\to\At$.}\;
Two rings appear: the matrix $V$ and its inverse live over
$\Z[q^{\pm1}]=\Z[v^{\pm2}]$, whereas the coordinates $c_r$ and the values
$\ev_m(g)$ live over $\Ft=\mathbb Q(v)[t^{\pm1}]$ --- honest $v$ (odd powers) may
occur, so the $c_r$ need not lie in $\mathbb Q(q)$. This causes no loss: the
extension $\Z[q^{\pm1}]\hookrightarrow\Z[v^{\pm1}]$ is free with basis
$\{1,v\}$, hence flat, and multiplication by a $\Z[q^{\pm1}]$-entry preserves
each of the two grades (even/odd powers of $v$). The reverse computation
$\mathbf c=V^{-1}\mathbf v$ therefore runs identically in each grade, so
integrality of the values over $\Z[v^{\pm1},t^{\pm1}]$ forces integrality of the
coordinates over the same ring. Equivalently, and most directly:
$V^{-1}$ has entries in $\Z[q^{\pm1}]\subseteq\At$ and $\mathbf v\in\At^{d+1}$,
so $\mathbf c=V^{-1}\mathbf v\in\At^{d+1}$.
\end{proof}

\subsubsection{Laurent clearing and conclusion}

\begin{proof}[Proof of Theorem~\ref{thm:iii}]
The inclusion $S\subseteq\Nev$ is Lemma~\ref{lem:fwd}. For $\Nev\subseteq S$,
let $f\in\Nev\subseteq\Ft[\Yc^{\pm1}]$ be arbitrary. Choose $N\ge0$ so large that
$g:=\Yc^{N}f\in\Ft[\Yc]$ is an honest polynomial (take $N$ at least the largest
magnitude of a negative $\Yc$-exponent occurring in $f$). For every $\chi\in\Z$,
using that $\ev_\chi$ is a ring homomorphism,
\[
  \ev_\chi(g)=\ev_\chi(\Yc^{N})\,\ev_\chi(f)=q^{N\chi}\,\ev_\chi(f).
\]
Here $q^{N\chi}$ is a unit of $\At$ and $\ev_\chi(f)\in\At$, so
$\ev_\chi(g)\in\At$; thus $g\in\Nev\cap\Ft[\Yc]$. By Lemma~\ref{lem:rev},
$g=\sum_{r}c_r\,\nu_r(\Yc)$ with $c_r\in\At$. Multiplying back by the unit
$\Yc^{-N}$,
\[
  f=\Yc^{-N}g=\sum_{r\ge0}c_r\,\Yc^{-N}\nu_r(\Yc)
   \ \in\ \sum_{u\in\Z}\sum_{r\ge0}\At\,\Yc^{\,u}\,\nu_r(\Yc)=S,
\]
each summand being an $\At$-multiple of the generator $\Yc^{u}\nu_r(\Yc)$ with
$u=-N$. Hence $\Nev\subseteq S$, and with Lemma~\ref{lem:fwd}, $\Nev=S$.
\end{proof}

\section{History}

\noindent
The lattice of Section~4 is the last term of a line that is additive before it is multiplicative
and classical before it is quantum.

\medskip\noindent
% NO YEAR ON THIS ONE, DELIBERATELY. The Library's Gramain card dates Polya's theorem 1914; the
% Attribution remark below - lifted VERBATIM from the checked proof and not mine to edit - dates it
% 1919. A paper cannot carry both, and I cannot settle which is right from what is held here, so
% this sentence carries no year rather than a plausible one. Reported to the Overseer.
\textbf{The classical case.} P\'olya's basis theorem: the polynomials taking integer values
at every integer form a free $\Z$-module on the binomial coefficients $\binom{x}{k}$, strictly
larger than $\Z[x]$ \cite[Prop.~1.1]{HarmanHopkins}. The subject it opened is surveyed in
Cahen--Chabert \cite{CahenChabert}.

\medskip\noindent
\textbf{The multiplicative case.} Replacing the integers by a geometric progression $1,q,q^2,\dots$
gives the grid that appears here. For entire functions the analogue is Gel'fond's (1933), which
Gramain records as the \emph{analogue multiplicatif} of P\'olya's theorem; for polynomials it is
Gramain \cite[Prop.~2.2]{Gramain}, who exhibits a basis for the functions taking integral values at
every $q^{n}$. One difference of setting matters: there $q$ is a natural number $\ge 2$ and
integrality means values in $\mathbb Z$, whereas here $q$ is an indeterminate and integrality means
values in $\At$. The grid is the same; the coefficient ring is not, and every structural
consequence drawn here from the free parameter $t$ is new with it.

\medskip\noindent
\textbf{The $q$-analogue.} Harman--Hopkins \cite[\S1, Props.~1.1--1.2]{HarmanHopkins} deform the
classical case: one variable, the additive grid $[n]_q$, coefficients in $\mathbb Z[q^{\pm1}]$, no
family parameter. The forward inclusion and the windowed reverse interpolation of Section~4 follow
their~\S1. What is added here is the passage to the multiplicative grid $\Yc=q^{\chi}$ and to
bilateral exponents $u\in\mathbb Z$, together with the Laurent clearing that derives the two-sided
statement from the one-sided one. Their substitution $\Yc=1+(q-1)x$ --- which their \S4 connects to
the Cartan part of Lusztig's integral form --- carries their grid node $x=[n]_q$ to $\Yc=q^{n}$ and
their basis term to $\nu_k(\Yc)$ term by term.

\medskip\noindent
\textbf{The quantum side.} For the $GL_n$ statement of Section~2 the ancestors are Lusztig
\cite[\S3.2]{Lusztig}, De Concini--Procesi \cite[\S\S17--18]{DCP}, Jantzen \cite[\S6]{Jantzen} and
Habiro--L\^e \cite[\S8G]{HabiroLe}; clause~(ii) rests on the corresponding $GL_2$ statement proved
earlier in this programme \cite{SandboxGL2}. Over the fraction field the centre is the invariant
\emph{monomial} ring. None of these treats the family parameter or the integer-valued lattice, and
the content of the flagship is precisely that over $\At$ the centre is the strictly larger Newton
lattice.

\begin{remark}
Clause~(iii) is the quantum P\'olya theorem of Harman and Hopkins,
\emph{Quantum integer-valued polynomials} (arXiv:1601.06110), in bilateral
Laurent form. Their \S1 Proposition (Prop.~\S1.2, \texttt{prop:qbinom}), the
$q$-analogue of P\'olya's 1919 theorem, states that the ring
$\mathcal R_q^{+}=\{P\in\mathbb Q(q)[x]:P([n]_q)\in\Z[q]\ \forall n\in\mathbb N\}$
is freely generated as a $\Z[q]$-module by the $q$-binomial coefficient
polynomials $\begin{bmatrix}x\\k\end{bmatrix}_{q}$ (the $q$-deformed Newton
basis on the \emph{additive} grid $x=[n]_q$, with denominator
$q^{\binom{k}{2}}[k]_q!$ --- not the classical P\'olya $\binom{x}{k}$); their \S4
Proposition (Prop.~\S4.3, \texttt{prop:localization}),
$\mathcal R_q=\mathcal R_q^{+}\otimes_{\Z[q]}\Z[q^{\pm1}]$, is the bilateral
localization. In the current notation their parameter $q$ coincides with our
$q$ ($=v^{2}$), so their displayed formulas read verbatim on our side.

The node-preserving affine substitution $\Yc=1+(q-1)x$ is exactly their \S4
$z$-variable; it carries the grid node $x=[n]_q$ to $\Yc=q^{n}$ and carries each
basis term $\begin{bmatrix}x\\k\end{bmatrix}_{q}$ onto $\nu_k(\Yc)$ \emph{term by
term} --- the factors $(q-1)$ cancel in every Newton ratio. (They remark further
that on adjoining square roots $K^{2}=z$, $v^{2}=q$ --- which is precisely the
passage to our $T^{+}$ and our $v$ --- the ring $\mathcal R_q[K,v]$ is equivalent
to the Cartan part of Lusztig's integral form of $U_v(\mathfrak{sl}_2)$. Nothing
below rests on that remark.) The values-in-$\Z[q^{\pm1}]$
side matches values-in-$\At$ with $t$ a spectator, since evaluation is
$t$-coefficient-wise and $\mathbb Q(q)\cap\Z[v^{\pm1}]=\Z[q^{\pm1}]$.

The two-variable ($GL_2$) form of this equivalence is established in
\cite{SandboxGL2}: there the grid criterion agrees with membership in the
$\At$-span of the product Newton generators, and passes to the full Laurent
even lattice. The $GL_1$ statement above is the rank-one specialization of
that transport.
\end{remark}

\begin{thebibliography}{9}

\bibitem{CahenChabert}
P.-J.~Cahen and J.-L.~Chabert,
\emph{Integer-Valued Polynomials},
Mathematical Surveys and Monographs \textbf{48}, American Mathematical Society, 1997.

\bibitem{DCP}
C.~De Concini and C.~Procesi,
\emph{Quantum groups},
in: D-modules, Representation Theory, and Quantum Groups,
Lecture Notes in Mathematics \textbf{1565}, Springer, 1993.

\bibitem{Gramain}
F.~Gramain,
\emph{Fonctions enti\`eres d'une ou plusieurs variables complexes prenant des valeurs enti\`eres
sur une progression g\'eom\'etrique},
in: Cinquante Ans de Polyn\^omes / Fifty Years of Polynomials,
Lecture Notes in Mathematics \textbf{1415}, Springer, 1990, pp.~123--137.

\bibitem{HabiroLe}
K.~Habiro and T.~T.~Q.~L\^e,
\emph{Unified quantum invariants for integral homology spheres associated with simple Lie
algebras},
Geom.\ Topol.\ \textbf{20} (2016), 2687--2835.

\bibitem{HarmanHopkins}
N.~Harman and S.~Hopkins,
\emph{Quantum integer-valued polynomials},
arXiv:1601.06110 (2016).

\bibitem{Jantzen}
J.~C.~Jantzen,
\emph{Lectures on Quantum Groups},
Graduate Studies in Mathematics \textbf{6}, American Mathematical Society, 1996.

\bibitem{Lusztig}
G.~Lusztig,
\emph{Introduction to Quantum Groups},
Progress in Mathematics \textbf{110}, Birkh\"auser, 1993.

\bibitem{SandboxGL2}
Project Sandbox,
\emph{The centre of the even hybrid family quantum $GL_2$}
(closure record of the signed goal, revision A1+R5), 2026.

\end{thebibliography}

\end{document}
